\chapter{Teoretická východiska a škálovací zákony}
\label{chap:theory}

\section{Architektura Decoder-Only Transformer}
V souladu s moderními autoregresivními jazykovými modely \cite{vaswani2017attention} využíváme architekturu Transformeru s kauzální maskou pozornosti, normalizací RMSNorm a aktivační funkcí SwiGLU.

\section{Škálovací zákony pro neuronové jazykové modely}
Kaplan et al. (2020) \cite{kaplan2020scaling} a Hoffmann et al. (2022) \cite{hoffmann2022training} popsali škálovací chování ztráty modelu $L$ jako funkci počtu ne-embeddingových parametrů $N$ a počtu trénovacích tokenů $D$:
\begin{equation}
    L(N, D) = E + \frac{A}{N^\alpha} + \frac{B}{D^\beta}
\end{equation}
kde $E$ představuje Bayesovskou ireducibilní entropii jazyka a $\alpha, \beta$ jsou škálovací exponenty.

\section{Metrika Bits-Per-Character (BPC)}
Standardní křížová entropie závisí na slovníku tokenizéru. Pro srovnání různých tokenizérů definujeme BPC vztažené k délce původního textu:
\begin{equation}
    \text{BPC} = \frac{\sum_{i} \mathcal{L}_i}{N_{\text{chars}} \cdot \ln 2}
\end{equation}
Tímto způsobem získáváme jednotnou škálu pro objektivní srovnání.
